% Replacement fragment for the HRES2--H2 part of paper/results.tex.
% This file is separate and does not modify the original manuscript.

\subsection{Results in the HRES2--H2 WPEB Engineering Case}
\label{sec:results_hres2_revision}

The engineering experiment evaluates the sizing and rule-based dispatch of a
grid-connected wind--photovoltaic--electrolyzer--battery (WPEB) system over
8,760 hourly steps. The 31 repetitions are stochastic optimization runs on a
common fixed synthetic weather profile; they should not be described as a
multi-year meteorological Monte Carlo experiment. The comparison comprises
the ten standalone methods PSO, GWO, WOA, EHO, ACO, ABC, ILS, SA, TS, and VNS.
The 31 corresponding executions were treated as paired runs for the inferential comparison.

The current implementation computes annual hydrogen production from the
electrolyzer load. It does not include a minimum hydrogen-demand constraint, a
hydrogen storage tank, or a fuel-cell subsystem. Accordingly, this result is a
WPEB capacity-sizing and dispatch case with hydrogen production, not a complete
hydrogen-storage microgrid.

\begin{table}[htbp]
\centering
\caption{Techno-economic performance metrics for the HRES2-H2 microgrid
($N=31$ stochastic optimization runs with corresponding executions treated as paired; 8,760-hour annual
simulation budget). Bold font indicates the best performance for each metric.}
\label{tab:hres2_metrics_revision}
\resizebox{\textwidth}{!}{
\begin{tabular}{lcccccc}
\toprule
\textbf{Algorithm} & \textbf{Mean LCOE} & \textbf{Best LCOE} & \textbf{Std. Dev. ($\sigma$)} & \textbf{LCOH} & \textbf{AGSR} & \textbf{Feasibility} \\
 & \textbf{(CNY/kWh)} & \textbf{(CNY/kWh)} & \textbf{(CNY/kWh)} & \textbf{(CNY/kg)} & \textbf{(\%)} & \textbf{Rate (\%)} \\
\midrule
\textbf{Proposed (DTW)} & \textbf{0.267160} & \textbf{0.267159} & \textbf{0.000001} & 17.6314 & \textbf{20.00\%} & \textbf{100.0\%} \\
\textbf{Proposed (DDTW)} & \textbf{0.267160} & \textbf{0.267159} & 0.000005 & \textbf{17.6313} & \textbf{20.00\%} & \textbf{100.0\%} \\
GWO & 0.268066 & 0.267160 & 0.002815 & 17.6800 & 19.42\% & 90.3\% \\
VNS & 0.267806 & 0.267159 & 0.002465 & 17.6700 & 19.95\% & 93.5\% \\
TS & 0.267188 & 0.267159 & 0.000038 & 17.6334 & 19.99\% & 96.8\% \\
ACO & 0.273288 & 0.267159 & 0.007584 & 18.0400 & 18.95\% & 93.5\% \\
ILS & 0.267184 & 0.267171 & 0.000023 & 17.6331 & 19.98\% & 96.8\% \\
PSO & 0.277289 & 0.274579 & 0.010939 & 18.2900 & 19.85\% & 83.8\% \\
WOA & 0.277227 & 0.276524 & 0.009567 & 18.3100 & 19.70\% & 87.1\% \\
EHO & 0.277728 & 0.276524 & 0.007382 & 18.3400 & 19.92\% & 80.6\% \\
SA & 0.286431 & 0.281339 & 0.013206 & 18.9100 & 19.35\% & 87.1\% \\
\bottomrule
\end{tabular}
}
\end{table}

\begin{table}[htbp]
\centering
\caption{Optimal physical system component sizing and annual green hydrogen
production for the HRES2-H2 microgrid. Bold font indicates the most
cost-effective component configuration.}
\label{tab:hres2_components_revision}
\resizebox{\textwidth}{!}{
\begin{tabular}{lccccc}
\toprule
\textbf{Algorithm} & \textbf{Wind Power} & \textbf{Solar PV} & \textbf{Electrolyzer} & \textbf{Battery Storage} & \textbf{Annual $\text{H}_2$ Production} \\
 & \textbf{(MW)} & \textbf{(MW)} & \textbf{(MW)} & \textbf{(MW / 4h)} & \textbf{(kg/year)} \\
\midrule
\textbf{Proposed (DTW)} & \textbf{174.47} & \textbf{25.53} & \textbf{70.0} & \textbf{50.0} & \textbf{8,862,879.4} \\
\textbf{Proposed (DDTW)} & \textbf{174.47} & \textbf{25.53} & \textbf{70.0} & \textbf{50.0} & \textbf{8,862,879.4} \\
GWO & 176.50 & 28.00 & 75.0 & 55.0 & 8,820,410.0 \\
VNS & 175.20 & 25.80 & 70.0 & 50.0 & 8,845,120.0 \\
TS & 174.50 & 25.60 & 70.0 & 50.0 & 8,860,250.0 \\
ACO & 178.00 & 27.50 & 75.0 & 55.0 & 8,790,500.0 \\
ILS & 175.00 & 26.00 & 70.0 & 50.0 & 8,855,100.0 \\
PSO & 185.00 & 35.00 & 85.0 & 65.0 & 8,650,200.0 \\
WOA & 182.00 & 32.00 & 80.0 & 60.0 & 8,690,400.0 \\
EHO & 190.00 & 40.00 & 90.0 & 70.0 & 8,610,000.0 \\
SA & 195.00 & 42.00 & 95.0 & 75.0 & 8,510,300.0 \\
\bottomrule
\end{tabular}
}
\end{table}

Both variants produced a stable low-LCOE solution on the archived scenario.
The best recorded configuration in both variants was
$P_{WT}=174.473154$ MW, $P_{PV}=25.526846$ MW, a $70$ MW electrolyzer, and a
$50$ MW battery with four hours of duration. This configuration produced
approximately $8.862879\times10^{6}$ kg of hydrogen per year, with LCOE
$0.267159$ CNY/kWh and LCOH $17.631389$ CNY/kg. The AGSR value was at or very
close to its $20\%$ upper bound, so the result should not be interpreted as
having a large constraint margin.

The archived LCOE statistics show significant differences among the eleven
methods according to the Friedman test for both variants. The proposed DTW and
DDTW variants obtained mean ranks of 1.71 and 1.95, respectively. All ten
standalone methods had higher mean LCOE ranks and unadjusted Wilcoxon
$p<0.001$ against the corresponding proposed control. The full inferential
comparison is reported in the companion statistical-analysis subsection; the
reported $p$-values are unadjusted because no Holm correction is applied.

The direct DTW--DDTW comparison yielded no significant difference in LCOE
(Wilcoxon $p=0.593$). The mean difference was only
$7.42\times10^{-7}$ CNY/kWh. Therefore, the engineering result supports
comparable performance of the two alignment modalities on the fixed scenario,
rather than a statistically established advantage of one modality.

The component table is retained as a descriptive summary of the reported
best configurations. However, the current standalone output does not preserve
per-run LCOH, AGSR, feasibility, hydrogen-production, or component-sizing
records. Therefore, the baseline engineering values in that table should be
interpreted as reported summary values rather than independently auditable
per-run comparisons until the complete records are archived.

\subsubsection{Stagnation-monitor dynamics}
\label{sec:results_hres2_dynamics_revision}

The convergence and DTW/DDTW profiles illustrate the switching behaviour of
the adaptive monitor over the 1,000-iteration budget. They should be described
as trajectories on the fixed synthetic engineering scenario. The final value
should be called the best observed solution rather than the global optimum,
because the metaheuristic search does not provide a proof of global
optimality.

The figure captions should also state ``ten standalone baseline
metaheuristics'' rather than ``nine standalone baseline metaheuristics'' and
should identify the experiment as 31 stochastic optimization runs over one
fixed synthetic weather profile.

% Implementation-aligned correction for the companion methods section:
% the photovoltaic temperature coefficient used by the code is
% gamma_T=-0.50 percent per degree Celsius, not -0.45 percent per degree.
